% Selected phase6 HH pairs only. No arbitrary-HC or uniform-family claim.
\paragraph{Two square support principles.}
Let $C,D$ be the two specified restored cycles and let $P,Q$ be
their pools. A legal Cartesian square $S$ has four edges.
If neither pool meets $E(S)$ and $C,D$ retain at least $a,b$
of those edges, respectively, then at least $a+b-4$ edges of $S$
belong to both protected sets. More generally, if both cycles
omit the two square edges at a common corner, the same bound is
$a+b-2$. These statements are just inclusion--exclusion in
$E(S)$, or in the two edges remaining after deleting that corner's
incident edges. They do not require the two retained edge sets
to be enumerated together.

For the six selected interfaces below, both pools avoid the
displayed square. Coordinates are numbered $0,\ldots,4$, and a
row gives its four words by independent choices of $x,y$.
The target uses the manuscript's first-normalizer convention.
\[
\begin{array}{c|c|c|c|c|c}
\text{source}&\text{change}&\text{square words}&x&y&(a,b)\\\hline
00014&0\to3&23xy1&1,2&0,3&(3,2)\\
01014&0\to q&xy r0q&1,r&0,q&(3,3)\text{ if }q=2;\ (2,3)\text{ if }q=3\\
01014&2\to2&xy320&1,3&0,2&(3,3)\\
01124&2\to2&x20y2&1,3&1,3&(2,3)\\
01204&2\to0&xy132&1,3&0,2&(3,2)
\end{array}
\qquad q\in\{2,3\},\quad r=5-q.
\]
All words are legal for both boundaries: each variable differs
from the fixed colours at its neighbours and from either
forbidden boundary colour. The listed cyclic words give the
retention counts independently on the two sides. Thus each
interface has a common protected edge; the two $(3,3)$ entries
have at least two. This is a six-case local support certificate,
not a formula for arbitrary selected cycles.

\paragraph{The remaining selected reflection pair.}
For $00123\to30123$, the prescribed normalizer restores the
same listed cycle by the involution
\[
g(w)_i=3-w_{-i}\qquad(i\bmod5).
\]
Consider the two legal Cartesian squares
\[
S=\{2x2y0:x,y\in\{1,3\}\},\qquad
T=\{13x1y:x,y\in\{0,2\}\}.
\]
The displayed source cycle restricted to $S$ is the path
$(21210,23210,23230)$ together with the internally isolated
corner $21230$; restricted to $T$, it is the path
$(13010,13210,13212)$ together with the internally isolated
corner $13012$. Indeed these two three-word paths occur
consecutively, and the isolated corners' cycle neighbours are
$(31230,21231)$ and $(13032,12012)$, respectively, all outside
their square.
The involution exchanges the squares, sends $23210$ to $13210$
and $21230$ to $13012$, and hence exchanges their internal
path centres and isolated corners. Consequently the restored
target cycle also retains two edges of $S$ and omits both edges
at $21230$. Neither the source pool nor the restored target
pool meets $S$ (the source pool avoids both $S$ and $T$).
The second square support principle yields two common
protected edges, without an actual common-edge list.

These conclusions concern the seven fixed selected pairs and
their prescribed maps. Simultaneous coordinate/colour
relabeling preserves each certificate. Arbitrary choices of
target normalizer or arbitrary Hamilton cycles are not asserted.
